% Clase 13 --- La ley de mallas es que el potencial está bien definido
% (§4.1 y §4.2).
% El circuito simple que dibujó el docente, y al lado el potencial a lo largo
% del recorrido. El panel (b) es el que explica la ley: al dar la vuelta
% completa hay que volver al mismo valor de V, y eso es todo lo que dice
% "suma de las diferencias de potencial igual a cero".
\begin{center}
\begin{tikzpicture}[scale=0.95]

  % =============== (a) el circuito, con el recorrido marcado ===============
  \begin{scope}
    \begin{scope}[line width=0.9pt, color=figblue]
      \draw (0,0) to[battery1, l=$\varepsilon$] (0,2.4);
      \draw (0,2.4) -- (3.2,2.4);
      \draw (3.2,2.4) to[R, l=$R$] (3.2,0);
      \draw (3.2,0) -- (0,0);
    \end{scope}

    % puntos del recorrido
    \foreach \p/\n/\pos in {(0,0)/A/{below left}, (0,2.4)/B/{above left},
                            (3.2,2.4)/C/{above right}, (3.2,0)/D/{below right}} {
      \node[figcarga, fill=figgray, minimum size=4pt] at \p {};
      \node[figlbl, \pos] at \p {$\n$};
    }

    % sentido del recorrido
    \draw[figvecres, line width=0.9pt] (1.3,1.2) arc (200:-70:0.45);
    \node[figlbl, figred, anchor=north] at (1.75,0.5) {\footnotesize recorrido};

    \draw[figvecres] (1.2,2.72) -- (2.1,2.72);
    \node[figlbl, figred, anchor=south] at (1.65,2.76) {$I$};

    \node[figlbl, anchor=north, align=center] at (1.6,-1.55)
      {(a) circuito simple};
  \end{scope}

  % =============== (b) el potencial a lo largo del recorrido ===============
  \begin{scope}[xshift=7.4cm, yshift=0.1cm]
    \draw[figaxis, line width=0.7pt] (-0.3,0) -- (6.1,0);
    \draw[figaxis, line width=0.7pt] (-0.3,0) -- (-0.3,2.9);
    \node[figlbl, anchor=east] at (-0.4,2.6) {$V$};
    \node[figlblaux, anchor=north] at (2.7,-0.95) {recorrido de la malla};

    % perfil: sube eps en la batería, constante en los cables, baja RI en R
    \draw[figblue, line width=1.1pt]
      (0,0.4) -- (1.0,0.4) -- (1.0,2.3) -- (3.2,2.3) -- (3.2,0.4) -- (5.4,0.4);
    \draw[figgray, dashed, line width=0.4pt] (1.0,0) -- (1.0,2.45);
    \draw[figgray, dashed, line width=0.4pt] (3.2,0) -- (3.2,2.45);
    \draw[figaux] (0,0.4) -- (5.6,0.4);

    \node[figlblaux, anchor=north] at (0.15,-0.08) {$A$};
    \node[figlblaux, anchor=north] at (1.25,-0.08) {$B$};
    \node[figlblaux, anchor=north] at (2.95,-0.08) {$C$};
    \node[figlblaux, anchor=north] at (3.45,-0.08) {$D$};
    \node[figlblaux, anchor=north] at (5.35,-0.08) {$A$};

    % las dos contribuciones
    \draw[figvecaux, figred, line width=0.9pt,
          -{Latex[length=2.2mm,width=1.6mm]}] (0.62,0.45) -- (0.62,2.25);
    \node[figlbl, figred, anchor=east] at (0.57,1.35) {$+\varepsilon$};
    \draw[figvecaux, figblue, line width=0.9pt,
          -{Latex[length=2.2mm,width=1.6mm]}] (3.6,2.25) -- (3.6,0.45);
    \node[figlbl, figblue, anchor=west] at (3.7,1.35) {$-R\,I$};

    \node[figlbl, figred, anchor=west] at (4.2,1.9) {$\varepsilon - R\,I = 0$};

    \node[figlbl, anchor=north, align=center] at (2.7,-1.55)
      {(b) al volver a $A$ el potencial es el mismo};
  \end{scope}

\end{tikzpicture}\\[0.5em]
{\small\itshape El panel (b) es la ley de mallas dibujada. Recorrer la malla es
ir leyendo el potencial punto a punto: sube $\varepsilon$ al atravesar la
batería del borne $-$ al $+$, se mantiene a lo largo de cada cable ---un cable
ideal es óhmico con $R=0$, así que no puede haber caída en él, y todos sus
puntos están al mismo potencial--- y baja $R\,I$ al atravesar la resistencia a
favor de la corriente. La única exigencia es la del final: al volver al punto de
partida hay que reencontrar el mismo valor, porque el potencial es una función
del punto. De ahí sale $\varepsilon = R\,I$ sin ninguna física nueva. Vale la
pena notar de dónde sale el sentido de la corriente: el potencial es mayor en el
borne positivo, hacia allí van los electrones, y la corriente convencional
resulta al revés que ellos.}
\end{center}
